\section{Specialization, connectedness, and the tuple count}\label{sec:specialization}

We retain the auxiliary fields and marked genus-one curve constructed in
\Cref{prop:mixed-lift,prop:spherical-base}. In this section $m=|J|$, $h\geq1$,
and $n=p^s$. Choose $s$ so that
\begin{equation}\label{eq:specialization-threshold}
 s\geq\max\{1,h\},\qquad
 p^{s-h}>m\bigl(h+\log_p m\bigr).
\end{equation}
This is a sufficient threshold for \Cref{thm:inseparable-genus}.
The relatively ample effective divisor $D$ has degree $n$ on each fiber
and is disjoint from the marked sections after their splitting extension.

\begin{proposition}\label{prop:no-small-cover}
Suppose $O_p(J)=1$ and \eqref{eq:specialization-threshold} holds. There is no
finite extension $L/K_s$ with $[L:K_s]_p<p^h$ over which the marked curve
$X=\mathcal X_L$ admits a geometrically connected smooth projective
$J$-Galois cover $f:Y\to X$, with its $J$-action defined over $L$, branched
only at the marked points and with ramification index divisible by $p$ at
every marked point.
\end{proposition}

We will obtain the contradiction by producing a nontrivial normal
$p$-subgroup of $J$. To this end, we first extract residue curves from
one $J$-orbit of valuations of a supposed cover. A comparison of
sections determines their total degree. A second comparison, with
prescribed vanishing, combines with the genus estimates to force
equality of Euler characteristics. These two equalities yield a flat
model whose special fiber is the disjoint union of those residue curves.
Connectedness will force the orbit to have one member; the kernel of
the resulting $J$-action on its residue field will be the required
nontrivial normal $p$-subgroup.

Suppose, for the proof, that such a cover exists. Its smooth,
geometrically connected source $Y$ is geometrically integral: over an
algebraic closure, the irreducible components of a smooth curve are
disjoint, so connectedness leaves only one. Let $R$ be the valuation
ring of $L$, with maximal ideal $\mathfrak m_R$ and residue field $F$.
By \Cref{prop:spherical-base}, $L$ is spherically complete, its value group
is $\mathbb Q$, and
\[
 [F:F_0]_p=[L:K_s]_p<p^h.
\]
Put $L'=LK_b'$ and $S=FF'$. By \Cref{lem:residue-permutations}, $S$ is the
residue field of $L'$, and the permutation actions on the labels $\xi_i$
over $L'$ and $t_i$ over $S$ agree. Write $D_X$ and $D_F$ for the divisors
induced by $D$ on $X$ and $E_F$, respectively, and put $D_Y=f^*D_X$.
The field extension $L(Y)/L(X)$ is Galois of degree $m$, with group $J$:
these assertions hold after extension to an algebraic closure, and the
specified automorphisms are already defined over $L$.

Let $e_i$ be the geometric ramification index at $\xi_i$. Each $e_i$ is
divisible by $p$. Riemann--Hurwitz over an algebraic closure, together
with invariance of Euler characteristic under field extension, gives
\begin{equation}\label{eq:generic-hurwitz}
 \frac{\nu(Y)}m
 =\frac{-2\chi(Y)}m
 =\sum_{i=1}^n\left(1-\frac1{e_i}\right).
\end{equation}
Here the base has genus one, and the reduced geometric fiber above
$\xi_i$ consists of $m/e_i$ points. The characteristic-zero form of
Riemann--Hurwitz being used is included in \Cref{foundation:complex}.

\subsection{The Gauss valuation and reduction of sections}

\begin{lemma}\label{lem:gauss-section-rule}
There is a valuation $v_X$ on $L(X)$ extending that on $L$, with value
group $\mathbb Q$ and residue field $F(E_F)$, having the following
properties. For every sufficiently large $j$, the integral ball in
$H^0(X,\mathcal O_X(jD_X))$ reduces onto
$H^0(E_F,\mathcal O_{E_F}(jD_F))$. Reduction is compatible with inclusions
as $j$ increases and with multiplication. Moreover, for an integral
section $b$ in such a space,
\begin{equation}\label{eq:gauss-evaluation}
 \overline{b(\xi_i)}=\bar b(t_i)
 \qquad\text{in }S.
\end{equation}
In particular, vanishing at $\xi_i$ implies vanishing of its reduction
at $t_i$.
\end{lemma}

\begin{proof}
Work initially over the complete DVR $R_b$, and write $D_E=D|_E$.
For $j$ sufficiently large,
\Cref{foundation:cohomology} and the special-fiber sequence give
\[
 H^0(\mathcal X,\mathcal O(jD))\otimes_{R_b}F_0
 \simeq H^0(E,\mathcal O_E(jD_E)).
\]
Indeed Serre vanishing kills $H^1(\mathcal X,\mathcal O(jD))$, so the
sequence obtained by multiplication by a uniformizer gives surjectivity
onto the special sections. The module on the left before reduction is
finite and torsion-free, hence free over $R_b$; torsion-freeness follows
from flatness of $\mathcal X$ and invertibility of $\mathcal O(jD)$.
Choose a basis $b_{j,1},\ldots,b_{j,t_j}$ whose reductions are a basis
of the special section space. Flat base change supplies bases over $L$
on the generic side and over $F$ on the special side.

As $D_X$ is effective, regard these sections as rational functions.
For $b=\sum_\alpha c_\alpha b_{j,\alpha}$ put
\[
 v_X(b)=\min_\alpha v(c_\alpha).
\]
Integral coefficients give a special rational function with poles
bounded by $jD_F$, and this function is nonzero if the displayed minimum
is zero. The inclusions for $jD\leq j'D$ are integral maps on the original
model and reduce to injective inclusions of special section spaces.
Consequently the assigned value is independent of the sufficiently
large space containing $b$: scale the coefficients so that their minimum
is zero, and use the nonzero special image. The same argument proves
multiplicativity. After scaling two nonzero sections to value zero, their
reductions are nonzero, and their product is nonzero because $E_F$ is
integral. The strong triangle inequality follows directly from the
coefficient rule.

The union of these section spaces is a ring whose fraction field is
$L(X)$: sufficiently high ample multiples give a projective embedding,
and their section ratios generate the function field. Thus the rule
extends to a valuation on $L(X)$. Its values lie in $\mathbb Q$, and all
values in $\mathbb Q$ occur already on $L$. A fraction of value zero can
be written with numerator and denominator both of value zero, by a
common scalar rescaling. Its residue is the ratio of their nonzero
special reductions, hence lies in $F(E_F)$. Conversely the ratios of
high-degree special sections generate $F(E_F)$, and each such section
has an integral lift. This identifies the residue field.

Finally the basis functions are regular along the splitting sections
over $R_b'$, since those sections avoid $D$. Their values are integral
and reduce to their evaluations at $t_i$. The same holds after extension
of coefficients from $R_b$ to $R$ and evaluation over the valuation ring
of $L'$. Since its residue is $S$, this proves
\eqref{eq:gauss-evaluation}.
\end{proof}

Choose one prolongation of $v_X$ to $L(Y)$, and let $\mathcal V$ be its
orbit of distinct conjugates under $J$. For $w\in\mathcal V$ let $K(w)$
be its residue field. The single-place residue-degree bound in
\Cref{lem:simultaneous-residues} makes $K(w)/F(E_F)$ finite; the Gauss
valuation has value group $\mathbb Q$ and infinite residue field $F(E_F)$.
Set
\[
 m_w=[K(w):F(E_F)].
\]
Let $C_w$ be the normalization of $E_F$ in $K(w)$ and $D_w$ the pullback
of $D_F$. These are regular integral projective curves, with finite
maps to $E_F$, by \Cref{foundation:normalization}.
Throughout the next argument we use this single orbit; its size is not
assumed in advance to account for all the extension degree.

For $j\geq0$ put
\[
 V_j=H^0(Y,\mathcal O_Y(jD_Y)),\qquad
 v_{\mathcal V}(a)=\min_{w\in\mathcal V}w(a),\qquad
 A_j=\{a\in V_j:v_{\mathcal V}(a)\geq0\}.
\]
The minimum is a valuation norm on the $L$-vector space $V_j$.
All prolongation values lie in $\mathbb Q$: algebraicity makes the
extension of value groups torsion, while the base value group is
divisible. The orthogonal-basis assertion of \Cref{lem:orthogonal-bases}
therefore applies, and an orthogonal basis can be scaled to have value
zero. In this basis,
\begin{equation}\label{eq:section-lattice}
 A_j\simeq R^{\dim_L V_j},\qquad
 \{a\in V_j:v_{\mathcal V}(a)>0\}=\mathfrak m_R A_j.
\end{equation}
These assertions also hold on any $L$-linear subspace of $V_j$.
At this point reduction means the map of rational functions
\[
 A_j\longrightarrow\prod_{w\in\mathcal V}K(w),
 \qquad a\longmapsto(\bar a_w)_w.
\]
Its kernel is the positive ball $\mathfrak m_R A_j$, so its image
has dimension $\dim_L V_j$ over $F$. No model of $Y$ has been used.
The next lemma locates these residue-field elements in the section
spaces of the separately normalized curves $C_w$.

\begin{lemma}\label{lem:section-specialization}
For all sufficiently large $j$, simultaneous residue gives an injection
\begin{equation}\label{eq:section-reduction}
 A_j/\mathfrak m_R A_j
 \lhook\joinrel\longrightarrow
 \bigoplus_{w\in\mathcal V}
 H^0(C_w,\mathcal O_{C_w}(jD_w)).
\end{equation}
These maps respect products. If a section vanishes on the reduced
inverse images of a selected set of marked points, its reductions
vanish on the corresponding reduced inverse images on the $C_w$,
provided the selections are defined over the respective ground fields.
The vanishing assertion may be checked after splitting the marked
points over $L'$ and $S$.
\end{lemma}

\begin{proof}
Let $a\in A_j$. Form the polynomial using every element of $J$,
including any repeated conjugates:
\begin{equation}\label{eq:conjugate-section-polynomial}
 P_a(T)=\prod_{g\in J}(T-g(a))
       =T^m+c_1T^{m-1}+\cdots+c_m.
\end{equation}
The coefficients are $J$-invariant rational functions, hence belong to
$L(Y)^J=L(X)$. At a fixed $w\in\mathcal V$ every
$g(a)$ is integral, because $g$ permutes $\mathcal V$. Thus each $c_k$
has nonnegative Gauss value. Moreover
\[
 c_k\in H^0(X,\mathcal O_X(kjD_X)).
\]
To see the pole bound, each conjugate has poles bounded by $jD_Y$, so
the elementary symmetric function of degree $k$ has poles bounded by
$kjD_Y$. Since the coefficient is a rational function on $X$, its
pole bound on $X$ follows from that on its finite surjective pullback.

By \Cref{lem:gauss-section-rule}, $\bar c_k$ has poles bounded by
$kjD_F$. The residue $\bar a_w\in K(w)$ satisfies the reduction of
\eqref{eq:conjugate-section-polynomial}. Let $t$ be a local equation for
$jD_F$ at any point of $E_F$. Then $t\bar a_w$ satisfies the monic
equation whose coefficient of degree $m-k$ is $t^k\bar c_k$.
These coefficients are regular there. At every point of $C_w$ above
this neighborhood, normality therefore makes $t\bar a_w$ regular.
This proves the required pole bound and hence
\eqref{eq:section-reduction}; injectivity was established in
\eqref{eq:section-lattice}. Residue maps themselves respect multiplication.

For the last assertion, work at a selected marked point after splitting.
Every $g(a)$ vanishes at every point of its reduced inverse fiber,
because this fiber is $J$-invariant. The non-leading coefficients
$c_k$ consequently vanish at the marked point downstairs. There are
no poles there, since $D_X$ avoids all the marked points.
By \eqref{eq:gauss-evaluation}, every $\bar c_k$ vanishes at $t_i$.
The reduced monic equation at any point above $t_i$ is then
$\bar a_w^m=0$ in its residue field, so $\bar a_w$ vanishes there.
\end{proof}

\subsection{Residue degrees and the inseparable quotient}

For $j$ sufficiently large, the line-bundle Euler-characteristic formula
and ample vanishing turn \eqref{eq:section-reduction} into
\begin{equation}\label{eq:full-section-comparison}
 jmn+\chi(Y)
 \leq jn\sum_{w\in\mathcal V}m_w
       +\sum_{w\in\mathcal V}\chi(C_w).
\end{equation}
Here $\deg D_X=\deg D_F=n$, and finite pullback multiplies divisor
degree by the map degree. Letting $j$ grow gives
$m\leq\sum_wm_w$. On the other hand,
\Cref{lem:simultaneous-residues}, applied to the Gauss valuation on
$L(X)$ and these distinct prolongations to the degree-$m$ extension
$L(Y)$, gives the reverse inequality. Its hypotheses hold: the value
group is $\mathbb Q$ and the residue field $F(E_F)$ is infinite.
We have proved
\begin{equation}\label{eq:residue-degree-equality}
 \sum_{w\in\mathcal V}m_w=m.
\end{equation}

Let $\Delta_w\leq J$ be the stabilizer of $w$, and let $I_w$ be the
kernel of its action on $K(w)$. The orbit has $m/|\Delta_w|$ members,
all of the same residue degree. Therefore
\begin{equation}\label{eq:decomposition-residue-degree}
 m_w=|\Delta_w|.
\end{equation}
In particular this equality follows from the section comparison, before
any assertion about the degree of an inseparable quotient.

\begin{lemma}\label{lem:inseparable-residue-quotient}
The invariant field $M_w=K(w)^{\Delta_w}$ is purely inseparable over
$F(E_F)$, of degree
\begin{equation}\label{eq:inertia-inseparable-degree}
 d=[M_w:F(E_F)]=|I_w|.
\end{equation}
In particular $d$ is a power of $p$ and $d\leq m$. If $Z_w$ denotes the
normalization of $E_F$ in $M_w$, the map $C_w\to Z_w$ is generically
separable Galois of degree $m_w/d$, with group $\Delta_w/I_w$.
All these towers are isomorphic over $E_F$ as $w$ varies in $\mathcal V$.
\end{lemma}

\begin{proof}
Take $x\in K(w)^{\Delta_w}$. Simultaneous approximation from
\Cref{lem:simultaneous-residues} gives an element $a\in L(Y)$ integral
at every member of $\mathcal V$, with residue $x$ at $w$ and zero at
the other members. The coefficients of the conjugate polynomial
$P_a(T)$ in \eqref{eq:conjugate-section-polynomial} have nonnegative
Gauss value. At $w$, its reduction is
\begin{equation}\label{eq:invariant-residue-polynomial}
 \overline{P_a}(T)
 =T^{m-|\Delta_w|}(T-x)^{|\Delta_w|}
 \in F(E_F)[T].
\end{equation}
Indeed, elements of $\Delta_w$ act on the residue $x$ trivially, and
each other conjugate draws its residue from a different member of
the orbit, where the prescribed residue is zero.

Write $|\Delta_w|=p^v q$ with $p\nmid q$. In characteristic $p$,
\[
 (T-x)^{|\Delta_w|}=(T^{p^v}-x^{p^v})^q.
\]
The coefficient of $T^{m-p^v}$ in
\eqref{eq:invariant-residue-polynomial} is therefore
$-q x^{p^v}$. The scalar $q$ is nonzero in $F(E_F)$, so
$x^{p^v}\in F(E_F)$. This also covers $x=0$, for which the coefficient
is zero. Thus $M_w/F(E_F)$ is purely inseparable.

The effective automorphism group on $K(w)$ is $\Delta_w/I_w$.
Finite fixed-field theory gives
\[
 [K(w):M_w]=|\Delta_w/I_w|.
\]
Combining this with \eqref{eq:decomposition-residue-degree} and the
tower formula proves \eqref{eq:inertia-inseparable-degree}.
The assertions about the normalizations and their separable function
fields follow. Conjugation by $J$ supplies the asserted isomorphisms
between the towers, so $d$ is independent of $w$.
\end{proof}

\subsection{Ramification on the residue curves}

Pass temporarily to the finite separable splitting field $S/F$.
The curve $Z_{w,S}$ remains integral: its finite map to $E_S$ is
surjective and radicial, while separable scalar extension preserves
regularity and hence reducedness. The curve $C_{w,S}$ is regular and
is a disjoint union of integral components. Each component dominates
$Z_{w,S}$: the finite flat map $C_{w,S}\to Z_{w,S}$ restricts on a
component to a nonzero finite locally free summand, of positive rank
on the connected target. The components are consequently the
normalizations in their respective function fields. These uses of
normalization and separable scalar extension are included in
\Cref{foundation:normalization,foundation:regularity}.

Let $z_i$ be the unique point of $Z_{w,S}$ above $t_i$, and write $r_i$
for its ramification index over $t_i$. Write $a_i$ for the ramification
index of $C_{w,S}\to Z_{w,S}$ above $z_i$. We justify that $a_i$ is
independent of the point above $z_i$, even when $C_{w,S}$ is disconnected.
Put $G_w=\Delta_w/I_w$ and $M_{w,S}=S(Z_{w,S})$. The generic algebra
of this map is
\[
 K(w)\otimes_{M_w}M_{w,S}=\prod_{\lambda=1}^t K_\lambda.
\]
It is a product of fields because $K(w)/M_w$ is finite separable.
The group $G_w$ acts on the first tensor factor, and the invariant
algebra is $M_{w,S}$: taking invariants is taking a kernel of linear
maps, which commutes with this scalar extension. An orbit of the
field factors gives an invariant idempotent. Since $M_{w,S}$ is a
field, there can be only one orbit.

The total dimension over $M_{w,S}$ is $|G_w|$, and transitivity makes
all factor degrees equal to $|G_w|/t$. The stabilizer of a factor has
this same order. It acts faithfully on that factor, since the map
from $K(w)$ into the factor is injective. Thus the factor extension
is Galois with that stabilizer as its group. Transitivity on primes
in each Galois normalization gives transitivity within each component
above $z_i$; the action on the factors joins these componentwise
orbits. We obtain transitivity on the entire fiber of $C_{w,S}$.
Ramification indices are preserved by this action. The isomorphisms
between the towers for different $w$ give the same $r_i,a_i$ throughout
the orbit.

Multiplicities in the tower multiply, so every point $P$ of $C_{w,S}$
above $t_i$ has index $r_i a_i$ over $t_i$. Finite flatness gives
\[
 \sum_{P\in C_{w,S},\ P\mid t_i}
    r_i a_i\,[\kappa(P):S]=m_w.
\]
Summing over $w$ and using \eqref{eq:residue-degree-equality} yields
\begin{equation}\label{eq:reduced-special-fiber-degree}
 \sum_{w\in\mathcal V}\ 
 \sum_{P\in C_{w,S},\ P\mid t_i}[\kappa(P):S]
 =\frac{m}{r_i a_i}.
\end{equation}
In particular all residue-field degrees, including inseparable ones,
are retained in the degree of the reduced fiber.

Apply \Cref{lem:separable-different} to $C_{w,S}\to Z_{w,S}$,
component by component. At a point above $z_i$, the determinant
vanishing is at least
$a_i-1+\mathbf 1_{p\mid a_i}$. The total separable degree of the
components over $Z_{w,S}$ is $m_w/d$. The preceding local-degree
calculation gives $m_w/(r_i a_i)$ for the sum of the full residue-field
degrees above $t_i$. Hence, for each $w$,
\[
 \nu(C_w)\geq\frac{m_w}{d}\nu(Z_w)
 +\sum_{i=1}^n\frac{m_w}{r_i a_i}
       \bigl(a_i-1+\mathbf 1_{p\mid a_i}\bigr).
\]
Euler characteristics and divisor degrees have the same numerical
values over $F$ as after the separable extension to $S$, so this
inequality is written in the original normalization over $F$.
The hypotheses of
\Cref{thm:inseparable-genus} hold for $Z_w/E_F$:
$[F:F_0]_p<p^h$, and \Cref{lem:inseparable-residue-quotient} gives
a purely inseparable degree $d\leq m$. Divide the displayed inequality
by $m_w$ and substitute that theorem's bound for $\nu(Z_w)/d$.
Now sum with weights $m_w/m$, whose sum is one by
\eqref{eq:residue-degree-equality}. This gives
\begin{equation}\label{eq:residue-genus}
 \frac{\sum_{w\in\mathcal V}\nu(C_w)}m
 \geq
 \sum_{i=1}^n\left(1-\frac1{r_i}\right)
 +\sum_{i=1}^n
   \frac{a_i-1+\mathbf 1_{p\mid a_i}}{r_i a_i}.
\end{equation}
Thus the normalization in this inequality agrees with that of
\eqref{eq:generic-hurwitz}.

\subsection{Vanishing conditions force equality}

The ramification estimate \eqref{eq:residue-genus} involves the
special indices $r_i a_i$, whereas \eqref{eq:generic-hurwitz}
involves the generic indices $e_i$. A second comparison of section
spaces will control their difference. Choose the labels in the set
\[
 \mathcal I=\left\{i:\frac1{r_i a_i}>\frac1{e_i}\right\}.
\]
At these labels, the degree $m/(r_i a_i)$ of the reduced special fiber
exceeds the degree $m/e_i$ of the reduced generic fiber. Imposing
vanishing there will make the section comparison detect this excess.
This selection descends on both sides. The generic index $e_i$ is
constant on the label orbits of $\operatorname{Gal}(L'/L)$, and the
special indices $r_i,a_i$ are constant on the label orbits of
$\operatorname{Gal}(S/F)$. These two permutation actions agree by
\Cref{lem:residue-permutations}. Hence the same subset of labels
defines a divisor over $L$ on the generic base and a divisor over $F$
on the special base.

Let $T_Y$ be the reduced inverse image on $Y$ of this selected divisor,
and let $T_w$ be its counterpart on $C_w$. These are effective divisors
on regular curves over their indicated fields. Their degrees may be
computed after the finite separable splitting extensions, which preserve
reducedness. The generic Galois cover and
\eqref{eq:reduced-special-fiber-degree} give
\begin{equation}\label{eq:selected-divisor-degrees}
 \deg_L T_Y=\sum_{i\in\mathcal I}\frac m{e_i},\qquad
 \sum_{w\in\mathcal V}\deg_F T_w
    =\sum_{i\in\mathcal I}\frac m{r_i a_i}.
\end{equation}

Restrict the norm to
$H^0(Y,\mathcal O_Y(jD_Y-T_Y))$. Its integral ball again has a
full-dimensional reduction by \eqref{eq:section-lattice}, and
\Cref{lem:section-specialization} places that reduction in
\[
 \bigoplus_{w\in\mathcal V}
 H^0(C_w,\mathcal O_{C_w}(jD_w-T_w)).
\]
For sufficiently large $j$, ample vanishing and
\eqref{eq:residue-degree-equality} therefore give
\[
 jmn+\chi(Y)-\deg_L T_Y
 \leq jmn+\sum_w\chi(C_w)-\sum_w\deg_F T_w.
\]
After rearranging and using \eqref{eq:selected-divisor-degrees}, this is
\begin{equation}\label{eq:vanishing-comparison}
 \frac{\nu(Y)}m
 \geq \frac{\sum_{w\in\mathcal V}\nu(C_w)}m
 +2\sum_{i=1}^n
   \max\left\{\frac1{r_i a_i}-\frac1{e_i},0\right\}.
\end{equation}

The three comparisons now force equality. To display their exact
interaction, put
\[
 q_i=\frac1{r_i a_i},\qquad b_i=\frac1{e_i},\qquad
 \delta_i=\mathbf 1_{p\mid a_i}.
\]
The contribution of label $i$ in \eqref{eq:residue-genus} is
$1-q_i+\delta_iq_i$. Combining that inequality with
\eqref{eq:vanishing-comparison} and subtracting
\eqref{eq:generic-hurwitz} gives
\begin{align}
 0
 &\geq \sum_{i=1}^n
       \bigl(b_i-q_i+2\max\{q_i-b_i,0\}+\delta_iq_i\bigr)
       \notag\\
 &=\sum_{i=1}^n\bigl(|q_i-b_i|+\delta_iq_i\bigr)
 \geq0.\label{eq:genus-nonnegative-excess}
\end{align}
Every summand vanishes, and every intervening inequality is an equality.
In particular,
\begin{equation}\label{eq:sharp-specialization}
 r_i a_i=e_i,\qquad p\nmid a_i\quad(1\leq i\leq n),
 \qquad d=|I_w|>1,
 \qquad \chi(Y)=\sum_{w\in\mathcal V}\chi(C_w).
\end{equation}
For the strict inequality $d>1$, use $p\mid e_i$ and $p\nmid a_i$:
then $r_i>1$, and $r_i$ divides the purely inseparable degree $d$ by
finite flatness at $t_i$. Also $q_i=b_i$ makes the extra term in
\eqref{eq:vanishing-comparison} zero, so equality there gives precisely
the displayed Euler-characteristic equality.

\subsection{A finitely presented model over the valuation ring}

\begin{proposition}\label{prop:specialization-model}
The degree equality \eqref{eq:residue-degree-equality} and the
Euler-characteristic equality in \eqref{eq:sharp-specialization}
produce a flat projective scheme $\mathcal Y$ of finite presentation
over $R$, with
\[
 \mathcal Y_L\simeq Y,\qquad
 \mathcal Y_F\simeq\coprod_{w\in\mathcal V}C_w.
\]
\end{proposition}

\begin{proof}
Form the graded $R$-algebra
\[
 A=\bigoplus_{j\geq0}A_j.
\]
Multiplication is defined because each $w$ is a valuation:
$A_iA_j\subseteq A_{i+j}$. Each piece is finite free by
\eqref{eq:section-lattice}. Moreover $A_0=R$, since $Y$ is projective
and geometrically integral and hence $H^0(Y,\mathcal O_Y)=L$.
Put
\[
 C=\coprod_{w\in\mathcal V}C_w,
 \qquad B_j=\bigoplus_{w\in\mathcal V}
 H^0(C_w,\mathcal O_{C_w}(jD_w)),
 \qquad B=\bigoplus_{j\geq0}B_j.
\]
The two equalities in the proposition make the source and target
dimensions in \eqref{eq:section-reduction} equal for every sufficiently
large $j$. Thus there is an integer $j_0$ such that
\begin{equation}\label{eq:full-high-degree-reduction}
 A_j/\mathfrak m_R A_j\simeq B_j
 \qquad(j\geq j_0),
\end{equation}
compatibly with products. Equality in each high degree, rather than
only equality of leading terms, is where the Euler-characteristic
equality is used.

We first choose a Veronese regrading for which the reduced algebra is
generated in degree one. The divisor on $C$ given by the $D_w$ is ample,
so $B$ is finitely generated by \Cref{foundation:cohomology}. Choose
positive-degree homogeneous generators $b_1,\ldots,b_t$ over $B_0$,
of degrees $d_1,\ldots,d_t$. Let $W$ be a common multiple of these
positive degrees. In a monomial in the $b_i$, remove pure powers
$b_i^{W/d_i}$, each of weight $W$, until the remaining exponent of
$b_i$ is less than $W/d_i$. The weight of the remainder is less than
$tW$. Choose $N=qW\geq j_0$ with $q\geq t$.

If the original monomial has degree $kN$, its remainder has weight
$cW$ for an integer $0\leq c<t\leq q$, since the removed blocks
and the original total weight are multiples of $W$. There are
$kq-c$ removed blocks. Combine the remainder with $q-c$ of these
blocks to make one factor of weight $N$, and group the remaining
$(k-1)q$ blocks in groups of $q$. This factors the monomial into
$k$ factors of weight $N$. Any coefficient from $B_0$ is absorbed
into the first factor. It follows that
\begin{equation}\label{eq:special-veronese-generation}
 F\ \oplus\!\bigoplus_{k\geq1} B_{kN}
 \quad\text{is generated over $F$ by }B_N.
\end{equation}
This argument permits $B_0=H^0(C,\mathcal O_C)$ to be larger than
$F$, as can happen for a disconnected curve or a curve with extra
constants. All the positive-degree factors used in the argument
belong to the full spaces in \eqref{eq:full-high-degree-reduction}.

Let $A^{(N)}=\bigoplus_{k\geq0}A_{kN}$, with $A_{kN}$ assigned degree
$k$. Choose an $R$-basis of $A_N$, and map the variables of a
degree-one polynomial ring
\[
 P=R[x_1,\ldots,x_u]\longrightarrow A^{(N)}
\]
to that basis. In each degree the target is a finite free $R$-module.
By \eqref{eq:full-high-degree-reduction} and
\eqref{eq:special-veronese-generation}, the map is surjective after
reduction modulo $\mathfrak m_R$. Its degreewise cokernel is finite,
so Nakayama's lemma makes the map surjective in every degree.
This proves finite generation over $R$.

For finite presentation, let $K_k$ be the kernel of
$P_k\to A_{kN}$. This surjection splits as an $R$-module map because
$A_{kN}$ is free. Hence $K_k$ is a finite direct summand of the finite
free module $P_k$, and reduction gives the exact kernel of the
corresponding polynomial presentation over $F$. The special relation
ideal is a homogeneous ideal in the Noetherian ring
$F[x_1,\ldots,x_u]$, so it has finitely many homogeneous generators.
Lift them, in their respective degrees, to elements $g_1,\ldots,g_v$
of $\ker(P\to A^{(N)})$, and let $I$ be the ideal they generate.
For every $k$, the quotient $K_k/I_k$ is a finite $R$-module. Its
reduction is zero, since the reductions of the $g_i$ generate the
special relation ideal. Degreewise Nakayama gives $K_k=I_k$ for every
$k$. Thus
\[
 A^{(N)}\simeq R[x_1,\ldots,x_u]/(g_1,\ldots,g_v)
\]
is a finite homogeneous presentation. In particular, this argument
does not require the nondiscrete valuation ring $R$ to be Noetherian.

Define $\mathcal Y=\operatorname{Proj}A^{(N)}$. The degree-one
presentation makes it a projective $R$-scheme of finite presentation.
It is flat over $R$: the graded algebra is a direct sum of free
$R$-modules, homogeneous localization is a filtered colimit of such
flat modules, and the degree-zero part of the resulting graded module
is a direct summand. Thus every standard affine chart of $\mathcal Y$
has a flat coordinate ring over $R$.

Since $A_j$ spans $V_j$, the generic fiber is the Proj of the ample
section ring of $Y$ after Veronese regrading, and hence is $Y$.
By \eqref{eq:full-high-degree-reduction}, the special fiber is the
Proj of the algebra in \eqref{eq:special-veronese-generation}.
Replacing its degree-zero term $F$ by $B_0$ has no effect on Proj.
Indeed, on a localization at a positive-degree element $f$, a missing
constant $b\in B_0$ is present as $(bf)/f$, since $bf$ has positive
degree. The degree-zero coordinate rings of these localizations
therefore agree. The full ample section ring recovers $C$, proving
the special-fiber assertion.
\end{proof}

\subsection{Connectedness and the final contradiction}

The finite presentation just proved is what permits passage to a
Noetherian base. This is the finite-equation approximation framework
of EGA IV \cite[Section~8]{EGAIV3}; the connected-fiber input is Stein
factorization, as in EGA III \cite[Section~4.3]{EGAIII1} and
\Cref{foundation:connectedness}. The following argument constructs
the normal Noetherian base explicitly, rather than applying coherent
finiteness or Stein factorization to the original valuation ring.

\begin{lemma}\label{lem:valuation-model-connectedness}
Let $R$ be a valuation ring of a characteristic-zero field $L$.
If a flat projective $R$-scheme $\mathcal U$ of finite presentation
has geometrically integral generic fiber, then its special fiber
is geometrically connected.
\end{lemma}

\begin{proof}
Choose a finite set of homogeneous projective equations for
$\mathcal U$. Their coefficients generate a finite-type
$\mathbb Z$-subdomain $R_0\subset R$. Normalize $R_0$ in its fraction
field, obtaining $R_1$. By \Cref{foundation:normalization}, this
normalization is finite, so $R_1$ is a normal Noetherian domain.
It remains a subring of $R$: its elements lie in $\operatorname{Frac}(R_0)$
and are integral over $R$, which is integrally closed.

The equations define a projective scheme $\mathcal U_1$ over $R_1$.
Its generic fiber is geometrically integral, because this can be
checked after extension of its fraction field to $L$, where the
fiber is $\mathcal U_L$. Let $\mathcal T\subseteq\mathcal U_1$ be the
scheme-theoretic closure of that generic fiber. Then $\mathcal T$ is
integral and projective over $R_1$: locally its coordinate ring is
the image in the coordinate ring of the integral generic fiber,
and the generic fiber is dense.

The pullback $\mathcal T_R$ is a closed subscheme of $\mathcal U$ and
agrees with it over $L$. Its defining ideal in
$\mathcal O_{\mathcal U}$ therefore vanishes upon localization at
$R\setminus\{0\}$. On each affine chart, every element of this ideal
is annihilated by a nonzero element of $R$. Flatness of $\mathcal U$
makes its affine coordinate rings torsion-free over $R$, so the ideal
is already zero. Thus
\begin{equation}\label{eq:normal-approximation-pullback}
 \mathcal T_R=\mathcal U.
\end{equation}

Write $K_1=\operatorname{Frac}(R_1)$ and
$\pi:\mathcal T\to\operatorname{Spec}R_1$.
Proper coherent finiteness over this Noetherian base gives a finite
$R_1$-algebra $H^0(\mathcal T,\mathcal O_{\mathcal T})$.
Restriction to the generic fiber is injective because $\mathcal T$
is integral and that fiber is dense. Since the generic fiber is
projective and geometrically integral, its global functions are
$K_1$. Hence
\[
 R_1\subseteq H^0(\mathcal T,\mathcal O_{\mathcal T})\subseteq K_1.
\]
The middle ring is finite over $R_1$, so all its elements are integral
over $R_1$. Normality yields
$H^0(\mathcal T,\mathcal O_{\mathcal T})=R_1$. Equivalently,
\[
 \pi_*\mathcal O_{\mathcal T}=\mathcal O_{\operatorname{Spec}R_1};
\]
one can check this equality on principal affine opens by flat base
change for coherent cohomology, as in \Cref{foundation:cohomology}.

Stein connectedness over the Noetherian base, in the form stated in
\Cref{foundation:connectedness}, now gives geometrically connected
fibers of $\pi$. The special fiber of $\mathcal U$ is, by
\eqref{eq:normal-approximation-pullback}, an extension of one of these
fibers to the residue field of $R$. It is therefore geometrically
connected. All uses of proper coherent cohomology and Stein
factorization in this proof occur over $R_1$.
\end{proof}

\begin{proof}[Completion of the proof of \Cref{prop:no-small-cover}]
Apply \Cref{lem:valuation-model-connectedness} to the flat projective
model of \Cref{prop:specialization-model}. Its generic fiber $Y$ is
geometrically integral, so its special fiber
$\coprod_{w\in\mathcal V}C_w$ is connected. Every $C_w$ is nonempty;
hence $\mathcal V$ has one member. The stabilizer of that member is
all of $J$, so $I_w$ is a normal subgroup of $J$. By
\eqref{eq:inertia-inseparable-degree} and
\eqref{eq:sharp-specialization}, it is a nonidentity $p$-group.
This contradicts $O_p(J)=1$.
\end{proof}

\begin{proof}[Proof of \Cref{thm:tuple-divisibility}]
Fix $h\geq1$ and choose $s_0$ satisfying
\eqref{eq:specialization-threshold} with $s=s_0$.
The same inequalities hold for every $s\geq s_0$; fix any such $s$.
For any prescribed multiset of
conjugacy classes of the $p$-singular puncture entries,
\Cref{lem:marked-cover-classification,prop:small-orbit-descent}
identify the finite tuple-class set with the corresponding marked-cover
classes and supply a Sylow
$p$-group action on that set. If an orbit had size less than $p^h$,
the same proposition would descend one of its covers to a finite
extension $L/K_s$ with $[L:K_s]_p<p^h$. This is excluded by
\Cref{prop:no-small-cover}.

Every orbit of a finite $p$-group has $p$-power size. Thus all the
orbits have sizes divisible by $p^h$, and so does their total.
The threshold depends only on $p,h,m$, not on the multiset.
An empty tuple set has count zero and satisfies the same assertion;
in particular this covers groups with no $p$-singular elements.
\end{proof}

\begin{proof}[Proof of \Cref{thm:main}]
The tuple divisibility just proved supplies the hypothesis of
\Cref{prop:group-reduction}. Therefore $l=(1+T)z$ for every pair
$(X,E)$ of a finite group and a central block sum considered there.
Taking $(X,E)=(G,B)$ gives
\[
 l(B)=
 \sum_{\substack{Q\leq G\text{ a }p\text{-subgroup}\\
                   \text{up to }G\text{-conjugacy}}}
 z\bigl(N_G(Q)/Q,\bar B_Q\bigr).
\]
The normalizer block-assignment calculation following
\Cref{lem:normalizer-projection} identifies this sum, including $Q=1$,
with $|W_p(B)|$ under the central-character convention
\eqref{eq:block-assignment}. Hence $l(B)=|W_p(B)|$ for every prime,
finite group, and block, as required.
\end{proof}
